\section{Resistência Ambiental e Cinética de Oxidação}\label{sec:resistencia_ambiental_oxidacao}

\subsection{O Desafio do Ambiente Hipersônico}\label{desafio-ambiente-hipersonico}

O voo hipersônico ($M > 5$) submete as estruturas a um fluxo aerotérmico severo, com temperaturas locais que variam de $1500~^\circ\mathrm{C}$ a mais de $2500~^\circ\mathrm{C}$, acompanhadas de espécies ionizadas de oxigênio e nitrogênio geradas pelo atrito de choque. Embora as Ligas de Alta Entropia Refratárias (RHEAs) mantenham rigidez mecânica superior em regimes térmicos elevados, a suscetibilidade à degradação termoquímica acelerada configura o principal obstáculo para a sua certificação aeroespacial.

\subsection{Mecanismos de Degradação e \emph{Pest Oxidation}}\label{mecanismos-degradacao-pest-oxidation}

A fragilidade ambiental das RHEAs decorre da presença majoritária de elementos refratários pesados (como W, Mo, Ta e Nb) com elevada afinidade termodinâmica pelo oxigênio. Em faixas intermediárias de temperatura, compreendidas entre $600~^\circ\mathrm{C}$ e $1000~^\circ\mathrm{C}$, manifesta-se o fenômeno destrutivo de \emph{pest oxidation}, no qual a peça sofre desintegração física e pulverização intergranular espontânea. 

Esse processo é impulsionado pelo descompasso volumétrico dos óxidos formados em relação à matriz metálica (fator de Pilling-Bedworth superior a 2), gerando tensões residuais compressivas que trincam e delaminam a camada de carepa superficial. Simultaneamente, a oxidação interna precipita compostos voláteis e de baixo ponto de fusão, como o trióxido de molibdênio ($\mathrm{MoO}_3$, que sublima a partir de $795~^\circ\mathrm{C}$) e óxidos de tungstênio ($\mathrm{WO}_3$). A volatilização gasosa dessas espécies abre caminhos de porosidade contínua, permitindo a difusão desimpedida de oxigênio gasoso em direção ao núcleo do material.

\subsection{Estratégias Composicionais e Passivação Superficial}\label{estrategias-composicionais-passivacao}

Para frear a cinética linear de degradação, a engenharia de materiais introduz elementos passivantes capazes de promover o crescimento de óxidos estáveis e protetores. A adição de alumínio estimula o desenvolvimento de uma camada contínua de alumina ($\alpha\text{-}\mathrm{Al}_2\mathrm{O}_3$), cuja baixa difusividade iônica reduz drasticamente as taxas de ganho de massa. Paralelamente, o cromo e o silício promovem subcamadas de cromia ($\mathrm{Cr}_2\mathrm{O}_3$) e sílica ($\mathrm{SiO}_2$), bloqueando a penetração de oxigênio mesmo em regimes dinâmicos. A dopagem com pequenas frações de terras raras, a exemplo do ítrio, ativa o conhecido \emph{pegging effect}, ancorando a carepa de óxido mecanicamente aos contornos da matriz e aumentando a vida útil sob fadiga térmica cíclica.

\subsection{Revestimentos Avançados e Hibridização Cerâmica}\label{revestimentos-avancados-hibridizacao}

Uma vez que a inclusão massiva de elementos passivantes pode comprometer a ductilidade à temperatura ambiente e induzir fases secundárias frágeis, o consenso atual aponta para a aplicação de revestimentos de proteção externa. Barreiras de silicietos modificados, como o $(\mathrm{NbMoTaW})\mathrm{Si}_2$ dopado com ítrio, operam com estabilidade até $2100~^\circ\mathrm{C}$. Sob calor extremo, esses revestimentos desenvolvem uma película viscosa de vidro de sílica que escorre e veda microtrincas de forma autorregenerativa (\emph{self-healing}). A integração desses sistemas a cerâmicas de ultra-alta temperatura (UHTCs) estabelece uma arquitetura híbrida sinérgica, em que a camada cerâmica externa dissipa a carga térmica enquanto o substrato refratário absorve os esforços mecânicos primários da estrutura.

\subsection{Evidências Cienciométricas na Literatura}\label{evidencias-cienciometricas-oxidacao}

O mapeamento quantitativo confirma que a termodinâmica de oxidação e a degradação ambiental constituem os temas de maior coesão intelectual na bibliografia recente. Redes de cocitação destacam que os grupos europeus de maior impacto dedicaram a última década a mapear diagramas de fases e cinética de recobrimento, consolidando o controle da \emph{pest oxidation} como critério indispensável para viabilizar as RHEAs em fuselagens e componentes hipersônicos operacionais.
